% Paso 4 (Alan). Tabla de requisitos del gimnasio.
% Mínimo 12 según la guía; la rúbrica pide 15, así que se hacen 15.
% Deben cubrir: negocio, stakeholder, usuario, funcional, NFR, interfaz, rendimiento, restricción y conformidad.
% RF en sintaxis EARS: "The [system] shall [function]" (en español: "El sistema deberá...").
% NFR: indicar la característica de la ISO/IEC 25010:2023.
% Interfaz: indicar el tipo (usuario, hardware, software, comunicaciones).
% Fuente: de dónde sale cada requisito (entrevista, formato real, norma...). Si aún no hay entrevista, se marca como supuesto.
\subsection{Tabla de tipos de requisitos aplicada al PFC}

{\small
\begin{xltabular}{\textwidth}{|>{\raggedright\arraybackslash}p{1.3cm}|>{\raggedright\arraybackslash}p{2.4cm}|L|>{\centering\arraybackslash}p{1.6cm}|>{\raggedright\arraybackslash}p{2.4cm}|}
\hline
\rowcolor{verdeuteq}\textcolor{white}{\textbf{ID}} & \textcolor{white}{\textbf{Tipo}} & \textcolor{white}{\textbf{Descripción}} & \textcolor{white}{\textbf{MoSCoW}} & \textcolor{white}{\textbf{Fuente}} \\ \hline
\endhead
RF-01 & Funcional & Cuando el encargado de turno registre la entrada de un estudiante, el sistema deberá guardar la fecha y la hora de entrada. & Must & Supuesto del equipo \\ \hline
RF-02 & Funcional & Cuando el encargado de turno registre la salida de un estudiante, el sistema deberá guardar la fecha y la hora de salida. & Must & Supuesto del equipo \\ \hline
RF-03 & Funcional & Cuando un estudiante no registre ninguna entrada durante 14 días hábiles consecutivos, el sistema deberá notificar al responsable del gimnasio cuál es ese estudiante. & Should & Supuesto del equipo \\ \hline
RF-04 & Funcional & El sistema deberá sugerir, mediante inteligencia artificial, una rutina para cada estudiante. & Could & Supuesto del equipo \\ \hline
RF-05 & Funcional & El sistema deberá asignar una rutina sugerida a un estudiante solo después de que el entrenador la apruebe. & Could & Supuesto del equipo \\ \hline
RF-06 & Funcional & El sistema deberá mostrar al responsable del gimnasio el número de estudiantes que tienen registrada una entrada sin su salida. & Must & Supuesto del equipo \\ \hline
RN-01 & Negocio & El sistema deberá permitir a Bienestar Universitario conocer qué estudiantes usan el gimnasio y en qué fechas y horas lo usan. & Should & Supuesto del equipo \\ \hline
RS-01 & Stakeholder & El encargado de turno necesita buscar a un estudiante por su número de cédula. & Must & Supuesto del equipo \\ \hline
RNF-01 & Rendimiento (ISO/IEC 25010: eficiencia de desempeño) & El sistema deberá mostrar el resultado de la búsqueda de un estudiante por cédula en un máximo de 3 segundos. & Must & Supuesto del equipo \\ \hline
RU-01 & Usuario & El estudiante necesita consultar la rutina que tiene asignada. & Could & Supuesto del equipo \\ \hline
RU-02 & Usuario & El estudiante necesita consultar sus registros de entrada y salida del gimnasio. & Must & Supuesto del equipo \\ \hline
RNF-02 & No funcional (ISO/IEC 25010: seguridad) & El sistema deberá permitir el acceso solo a usuarios que se autentiquen con usuario y contraseña. & Must & Supuesto del equipo \\ \hline
RI-01 & Interfaz (software) & El sistema deberá permitir importar el listado de estudiantes habilitados que entregue Bienestar Universitario cada período académico. & Should & Supuesto del equipo \\ \hline
RR-01 & Restricción & El sistema no deberá conectarse directamente con el sistema académico de la UTEQ; la habilitación de los estudiantes se validará con el listado de Bienestar Universitario. & Must & Supuesto del equipo \\ \hline
RC-01 & Conformidad & El sistema deberá tratar los datos personales de los estudiantes conforme a la Ley Orgánica de Protección de Datos Personales del Ecuador (2021). & Must & LOPDP (2021) \\ \hline
\end{xltabular}
}
